\section{Text Embeddings and Neural Retrieval: State of the Art}
\label{sec:embeddings}

Neural text retrieval rests on learned dense representations of text that can be
compared by an inexpensive geometric operation. This section surveys the
architectural families that dominate the field, the contrastive objectives that
train them, the recent wave of instruction-tuned and LLM-based embedders, and the
representation-geometry choices (dimensionality, pooling, Matryoshka nesting) that
govern their deployment. Throughout we tie the landscape back to the project's own
bi-encoder/cross-encoder findings on Adobe Experience Platform (AEP) and Adobe
Journey Optimizer (AJO) workflow data, and flag the one assumption---\emph{symmetric}
cosine similarity---that makes all of these models a poor fit for directional, causal
ordering.

\subsection{Architecture families}
\label{sec:emb:arch}

\paragraph{Bi-encoders (dual-encoders).}
A bi-encoder maps a query $q$ and a document $d$ \emph{independently} to vectors
$\mathbf{e}_q,\mathbf{e}_d\in\mathbb{R}^k$ and scores them with a symmetric kernel,
typically inner product or cosine: $s(q,d)=\mathbf{e}_q^{\!\top}\mathbf{e}_d$. Because
document vectors can be precomputed and indexed, retrieval reduces to (approximate)
nearest-neighbour search and scales to billions of passages. DPR
\citep{emb:karpukhin2020dpr} established the modern recipe of two BERT towers trained
with in-batch negatives; Sentence-BERT \citep{emb:reimers2019sbert} popularised
siamese/triplet fine-tuning for sentence-level similarity; and SimCSE
\citep{emb:gao2021simcse} showed that even unsupervised dropout-as-augmentation
contrastive training yields strong sentence embeddings while reducing the
anisotropy of raw BERT representations \citep{emb:li2020sentencebias}. The repo's
\texttt{DirectionalClassifier} is a bi-encoder variant: two untied encoders produce
CLS embeddings that are L2-normalised and fused as $[\mathbf{a};\mathbf{b};
\mathbf{a}{-}\mathbf{b};\mathbf{a}{\odot}\mathbf{b}]$ before an MLP head---the
explicit $\mathbf{a}{-}\mathbf{b}$ term is an attempt to inject the asymmetry that
plain cosine lacks.

\paragraph{Cross-encoders (rerankers).}
A cross-encoder concatenates both texts into a single input
\texttt{[CLS]\,$t_1$\,[SEP]\,$t_2$} and lets full self-attention model token-level
interactions, emitting a scalar relevance from the pooled state
\citep{emb:nogueira2019passage}; sequence-to-sequence variants such as monoT5
\citep{emb:nogueira2020monot5} cast ranking as generating a ``true''/``false'' token.
Cross-encoders are far more accurate but cost a full forward pass \emph{per pair},
so they are used only as second-stage rerankers over a shortlist from a bi-encoder.

\paragraph{Late interaction (ColBERT family).}
ColBERT \citep{emb:khattab2020colbert} occupies a middle ground: it stores
\emph{per-token} embeddings and scores a pair by a sum of max-similarities
(``MaxSim''), $s(q,d)=\sum_{i\in q}\max_{j\in d}\mathbf{e}_{q_i}^{\!\top}\mathbf{e}_{d_j}$,
preserving fine-grained interaction while keeping the document side precomputable.
ColBERTv2 \citep{emb:santhanam2022colbertv2} adds residual compression and
denoised, distilled supervision, cutting the index footprint $6$--$10\times$, and
PLAID \citep{emb:santhanam2022plaid} introduces centroid-based pruning that lowers
late-interaction latency by up to $7\times$ on GPU and $45\times$ on CPU at equal
quality.

\paragraph{The accuracy/latency tradeoff and the repo finding.}
The canonical wisdom---cross-encoders are more accurate, bi-encoders are faster---is
borne out and \emph{sharpened} by the project's experiments. On pairwise
\emph{classification} (\texttt{aep\_causal\_cls34}), the bi-encoder \emph{beat} the
cross-encoder by roughly $4.6\%$ (e.g.\ BGE-base $86.1\%$ bi vs.\ $81.5\%$ cross;
MiniLM $87.4\%$ vs.\ $78.8\%$), because the cross-encoder shares a single $512$-token
budget across both texts and is \emph{starved} of context when the two steps are long.
Yet on the task we actually care about---end-to-end \emph{workflow ordering}---the
DeBERTa-v3-large cross-encoder scored $21/30$ while the bi-encoder scored $0/30$. The
reason is distributional: training pairs are same-document and windowed (adjacent
steps, gap $\le 2$--$4$), but ordering a full workflow requires judging
\emph{non-adjacent, off-distribution} pairs. The bi-encoder's independent embeddings
collapse there---two far-apart steps look semantically similar and cosine cannot tell
which precedes---whereas the cross-encoder's \emph{joint cross-attention} reads the two
texts \emph{together} and can pick up directional cues (prerequisite mentions,
references to prior configuration) that survive the distribution shift. Pairwise
accuracy is therefore \emph{not} a proxy for ordering quality.

\subsection{Contrastive training and negatives}
\label{sec:emb:contrastive}

Dense encoders are trained almost universally with the InfoNCE / softmax-contrastive
loss \citep{emb:oord2018infonce}. For a query $q$ with one positive $d^+$ and
negatives $\{d_j^-\}$,
\[
\mathcal{L}_{\mathrm{InfoNCE}}=-\log
\frac{\exp\!\big(s(q,d^+)/\tau\big)}
{\exp\!\big(s(q,d^+)/\tau\big)+\sum_j \exp\!\big(s(q,d_j^-)/\tau\big)},
\]
where $\tau$ is a temperature controlling how sharply the model separates positives
from negatives (small $\tau$ emphasises the hardest negatives). \emph{In-batch
negatives} reuse other examples in the minibatch as cheap negatives, so larger
batches directly improve representation quality. The dominant lever, however, is
\emph{hard-negative mining}: negatives that are lexically or semantically close to the
positive but actually irrelevant. ANCE \citep{emb:xiong2021ance} periodically re-mines
hard negatives from the model's own evolving ANN index; RocketQA
\citep{emb:qu2021rocketqa} combines cross-batch negatives, denoised hard negatives
(filtered by a cross-encoder to remove false negatives), and data augmentation; and
TAS-B \citep{emb:hofstatter2021tasb} balances topic-aware sampling with
cross-encoder/ColBERT distillation. This directly motivates the repo's
\texttt{hard\_neg\_semantic} dataset ($\sim$261K pairs): forcing the model to
distinguish a true causal successor from a topically identical but causally wrong step
is exactly the hard-negative regime, and it is what pushes the encoder past the
``same-topic'' shortcut toward genuine directional discrimination.

\subsection{Instruction-tuned and LLM-based embeddings}
\label{sec:emb:llm}

The current frontier replaces BERT-scale encoders with instruction-conditioned and
decoder-LLM backbones. E5 \citep{emb:wang2022e5} pre-trains on weakly supervised text
pairs (CCPairs) and was the first to beat BM25 zero-shot on BEIR; its multilingual
variant underpins much of the open ecosystem. INSTRUCTOR \citep{emb:su2023instructor}
prepends a natural-language \emph{task instruction} so one model serves many tasks,
and GTE \citep{emb:li2023gte} uses multi-stage contrastive training over diverse
sources. The BGE family \citep{emb:xiao2023bge} and BGE-M3 \citep{emb:chen2024bgem3}
add self-knowledge distillation and, in M3, \emph{multi-functionality} (dense,
sparse, and multi-vector retrieval from one model), $100{+}$ languages, and
$8192$-token context. The LLM-backbone line is now state of the art: E5-Mistral
\citep{emb:wang2024mistrale5} fine-tunes Mistral-7B on GPT-generated synthetic
retrieval tasks; SFR-Embedding-Mistral \citep{emb:sfr2024embedding} extends it with
multi-task transfer; LLM2Vec \citep{emb:behnamghader2024llm2vec} converts any
decoder-only LLM into an encoder via bidirectional attention plus masked-next-token
and contrastive adaptation; NV-Embed \citep{emb:lee2025nvembed} adds latent-attention
pooling and two-stage instruction tuning to reach $72.31$ on MTEB; GritLM
\citep{emb:muennighoff2024gritlm} unifies generation and embedding in one model via
instruction routing; and Qwen3-Embedding \citep{emb:zhang2025qwen3emb} (0.6B/4B/8B)
tops the MMTEB multilingual leaderboard ($70.58$ for the 8B as of mid-2025), with even
the \emph{0.6B} model competitive---directly relevant under our $<\!1$B budget.

\paragraph{What the leaderboards measure---and don't.}
MTEB \citep{emb:muennighoff2023mteb} aggregates $\sim$56 tasks (retrieval/BEIR
\citep{emb:thakur2021beir}, clustering, classification, STS, reranking,
pair-classification, summarisation) into a single average; MMTEB
\citep{emb:enevoldsen2025mmteb} extends this to $500{+}$ tasks and $250{+}$ languages
with long-document, code, and instruction-following tasks. Crucially, \emph{none} of
these tasks is \emph{directional} or \emph{causal}: STS and pair-classification are
symmetric ($\mathrm{sim}(a,b)=\mathrm{sim}(b,a)$), retrieval is topical relevance, and
reranking judges relevance, not precedence. A model can top MMTEB and still be unable
to tell whether step $A$ \emph{enables} step $B$ or vice versa. This is borne out by
the project's own GritLM-7B probe: zero-shot, it \emph{stays semantic}---it ranks by
topical similarity and \emph{ignores} a causal instruction in the prompt, never
overriding its semantic bias in favour of directional precedence.

\subsection{Representation geometry: dimensionality, pooling, Matryoshka}
\label{sec:emb:geometry}

\paragraph{Pooling.} How token states become one vector matters. CLS pooling reads a
single special token (natural for BERT bi-encoders and the repo's classifier); mean
pooling averages token states (the default for E5/Sentence-BERT and usually more
robust); last-token pooling reads the final position and is standard for causal
decoder-LLM embedders (E5-Mistral, NV-Embed, Qwen3) because only that position has
attended to the whole sequence---NV-Embed's latent-attention pooling is an explicit
fix for the weakness of plain last-token pooling.

\paragraph{Dimensionality and Matryoshka.} Larger $k$ (LLM embedders emit
$1024$--$4096$-d) helps quality but inflates index size and search cost. Matryoshka
Representation Learning \citep{emb:kusupati2022matryoshka} trains nested prefixes so a
single vector can be \emph{truncated} (e.g.\ $1024\!\to\!256\!\to\!64$) with graceful
degradation, enabling an adaptive accuracy/latency knob and cheap coarse-to-fine
retrieval; it is now standard in commercial APIs and in BGE/Qwen3 releases. For a
two-stage AEP/AJO system, MRL allows a small truncated vector for stage-1 recall and
the full vector (or a reranker) for precision.

\paragraph{The symmetry mismatch.} Every model above scores pairs with cosine or inner
product, which is \emph{symmetric}: $s(a,b)=s(b,a)$. A retrieval relevance signal is
naturally symmetric, but causal precedence is \emph{antisymmetric}---if $A$ precedes
$B$ then $B$ does not precede $A$. No amount of contrastive fine-tuning of a symmetric
kernel can represent a strict order, which is precisely why the repo's bi-encoder bolts
on an $\mathbf{a}{-}\mathbf{b}$ feature and why the cross-encoder (asymmetric by
input order) wins on ordering. This is the central geometric gap addressed later (see
the section on asymmetric geometry), which considers order embeddings, asymmetric/dot
scoring, and hyperbolic/cone representations that natively encode direction.

\subsection*{Relevance to causal embedding for AEP/AJO}

Three implications follow for the project. (1) A symmetric-cosine bi-encoder is an
excellent \emph{stage-1 recaller}---fine-tuned BGE-large already gives $+21\%$ MRR over
the BERT baseline on AEP retrieval---but its independent, symmetric embeddings
\emph{cannot} order non-adjacent workflow steps, matching the $0/30$ ordering result.
(2) The cross-encoder's joint attention is what generalises to off-distribution pairs
($21/30$); the right architecture is therefore a bi-encoder retriever feeding an
\emph{asymmetric, cross-attentive} reranker/orderer. (3) Instruction-tuned LLM
embedders are tempting under the $<\!1$B budget (Qwen3-Embedding-0.6B is strong), but
the GritLM finding warns that instructions alone do not buy directionality---these
models optimise the symmetric, topical objective that MTEB/MMTEB reward, with no causal
or directional task in sight. Closing that gap requires hard causal negatives
(\S\ref{sec:emb:contrastive}), an antisymmetric scoring geometry, and cross-attention,
not merely a bigger embedder.
